\section{Continuation to the full spacetime tube}
\label{sec:full-tube}
The ordered complex-time functions constructed earlier are holomorphic in
the time variables. Locality needs a larger domain: the spatial variables
must also become complex. We construct this continuation from the spectral
support of the boundary distributions. A smooth cutoff near that support
makes the Fourier--Laplace expression an actual Schwartz test. Keeping the
cutoff also permits limits in which some imaginary gaps lie on the boundary
of the forward cone.

Fix $n$ and set $D=4n$. All coordinates in this section are absolute
coordinates in $\mathbb R^D$, rather than relative coordinates. Write
\begin{align*}
 K_n&=\left\{p:\sum_i p_i=0,\quad
                  \sum_{i=j}^n p_i\in\overline V_+\ (2\leq j\leq n)\right\},\\
 C_n&=\{y:y_{i+1}-y_i\in\overline V_+\ (1\leq i<n)\},\\
 C_n^\circ&=\{y:y_{i+1}-y_i\in V_+\ (1\leq i<n)\},
 \qquad V_+=\{q:q^0>|\mathbf q|\}.
\end{align*}
The superscript on $C_n^\circ$ denotes the strict-gap set, which is open
in the full space. For $n=0$ or $1$ the gap conditions are vacuous.
The tube is $\mathcal T_n=\mathbb R^D+iC_n^\circ$.
By Corollary~\ref{cor:cone-support} and
Theorem~\ref{thm:lorentz-covariance}, the distribution
$\widehat W_n(v)$ has support in $K_n$, under the hypotheses stated there.

\begin{lemma}[Dual positivity and a strict margin]\label{lem:tube-margin}
If $p\in K_n$ and $y\in C_n$, then $p\cdot y\geq0$.
For each $y\in C_n^\circ$ there are $c>0$ and a neighborhood $N$ of
$y$ such that
\[
 p\cdot u\geq c|p|\qquad(p\in K_n,\ u\in N).
\]
\end{lemma}
\begin{proof}
Put $q_j=\sum_{i=j}^n p_i$. The zero total momentum and telescoping give
\[
 p\cdot y=\sum_{j=2}^n q_j\cdot(y_j-y_{j-1}).
\]
The Euclidean pairing of two forward-cone vectors is nonnegative:
$a\cdot b=a^0b^0+\mathbf a\cdot\mathbf b
\geq a^0b^0-|\mathbf a||\mathbf b|\geq0$.
This proves the first assertion. Choose $r>0$ with
$B(y,r)\subseteq C_n^\circ$, and put $c=r/2$.
For $u\in B(y,r/2)$ and $p\ne0$, the point
$u-cp/|p|$ belongs to $B(y,r)$ and hence to $C_n$.
Applying the first assertion to this point gives
$p\cdot u-c|p|\geq0$. The case $p=0$ is immediate.
This argument also covers the vacuous gap cases, where $K_n=\{0\}$.
\end{proof}

\subsection{A smooth cutoff and closed imaginary shifts}
Let $K_n^{(a)}=\{p:\operatorname{dist}(p,K_n)<a\}$.
Choose a nonnegative smooth bump $\rho$, supported in $B(0,1/2)$,
with integral one, and define
\[
 \chi=\rho*\mathbf1_{K_n^{(2)}}.
\]
Then $0\leq\chi\leq1$, $\chi=1$ on $K_n^{(1)}$, and
$\operatorname{supp}\chi\subseteq K_n^{(3)}$, with support understood
as closed support. Indeed, for $p\in K_n^{(1)}$ every point $p-q$ with
$q\in\operatorname{supp}\rho$ has distance less than $3/2$ from $K_n$;
the convolution therefore integrates the constant one. Nonzero values
require distance less than $5/2$ from $K_n$, and their closure still has
distance at most $5/2<3$.
Every derivative is globally bounded, since
\[
 \partial^\alpha\chi=(\partial^\alpha\rho)*\mathbf1_{K_n^{(2)}},
 \qquad \|\partial^\alpha\chi\|_\infty
             \leq\|\partial^\alpha\rho\|_1.
\]
Differentiating the compactly supported convolution proves smoothness:
difference quotients of the smooth factor converge in $L^1$, so their
convolutions with the bounded indicator converge uniformly. In dimension
zero simply take $\chi=1$.

For $y\in C_n$ set $a_y(p)=e^{-y\cdot p}\chi(p)$.
This is a smooth multiplier of Schwartz space. More precisely, for every
$R$ and multi-index $\alpha$ there is a finite constant $A_{R,\alpha}$
such that
\begin{equation}\label{eq:closed-damper-bound}
 \sup_p|\partial^\alpha a_y(p)|\leq A_{R,\alpha}
       \qquad(y\in C_n,\ |y|\leq R).
\end{equation}
On the closed support of $\chi$ choose $q\in K_n$ with $|p-q|<3$.
Dual positivity gives $-y\cdot p\leq3|y|$. Each derivative of the
exponential adds a product of components of $y$, while derivatives of
$\chi$ are bounded. Leibniz's rule proves
\eqref{eq:closed-damper-bound}; outside that support all derivatives of
$a_y$ vanish.

For each $f\in\mathscr S$, the map $y\mapsto a_yf$ is continuous from
$C_n$ to $\mathscr S$. On a compact frequency ball this follows uniformly
for every derivative from the explicit smooth formula. Outside a large
ball, the uniform derivative bounds and Leibniz's rule bound a seminorm
of weight $k$ by a constant times $(1+|p|)^{-1}$ times finitely many
seminorms of $f$ of weight $k+1$. The tail is therefore uniformly small.
Combining the compact and tail estimates proves convergence in every
Schwartz seminorm, including at closed-cone boundary shifts.

For a fixed label tuple abbreviate $w=\widehat W_n(v)$ and define a
tempered distribution $W_y$ by
\begin{equation}\label{eq:partial-damping}
 W_y(f)=w(a_y\check f),\qquad
 \check f(p)=(2\pi)^{-D}\int e^{ip\cdot x}f(x)\,dx.
\end{equation}
The preceding continuity shows that $W_y(f)$ is continuous in
$y\in C_n$. Moreover $W_0=W_n(v)$. To see this, multiplication by
$\chi$ leaves $w$ unchanged because $\chi=1$ on a neighborhood of
its support. The support assertion applies to Schwartz tests as well as
compact tests: multiply a test vanishing near the support by compact
cutoffs, apply distributional support, and pass to the Schwartz limit.
Fourier inversion now gives $w(\check f)=W_n(v)(f)$.

\subsection{The holomorphic Fourier--Laplace expression}
For $z=x+iy\in\mathcal T_n$ define
\[
 L_z(p)=e^{ip\cdot z}\chi(p),\qquad
 H_n(v;z)=(2\pi)^{-D}w(L_z).
\]
Here $p\cdot z$ is the complex bilinear extension of the Euclidean
pairing, with no conjugation. The test $L_z$ is in $\mathscr S$.
In fact, the margin lemma at $y$ and the distance-three support property
give, on $\operatorname{supp}\chi$,
\[
 -y\cdot p\leq3(|y|+c)-c|p|.
\]
To obtain this, choose $q\in K_n$ with $|p-q|<3$, use
$y\cdot q\geq c|q|$, and $|q|\geq|p|-3$.
Frequency derivatives of $L_z$ only introduce bounded derivatives of
$\chi$ and powers of the fixed vector $z$. They consequently have
exponential decay. The same estimates hold in a sufficiently small
neighborhood of $z$, with a common positive decay rate.

\begin{proposition}\label{prop:full-tube-holomorphic}
The map $z\mapsto L_z$ is complex differentiable into Schwartz space,
with derivative $h\mapsto i(p\cdot h)L_z(p)$. Consequently
$H_n(v;z)$ is holomorphic jointly in all $D$ complex coordinates.
For each fixed $y\in C_n^\circ$ it grows at most polynomially in $x$,
and for every $f\in\mathscr S$,
\begin{equation}\label{eq:tube-integral}
 W_y(f)=\int H_n(v;x+iy)f(x)\,dx,
\end{equation}
with an absolutely convergent integral.
\end{proposition}
\begin{proof}
The difference from the proposed derivative is
\[
 L_{z+h}-L_z-i(p\cdot h)L_z
 =L_z\big(e^{ip\cdot h}-1-ip\cdot h\big).
\]
For $|h|\leq1$, every fixed-order frequency derivative of the factor
in parentheses has absolute value at most
$C_\alpha |h|^2(1+|p|)^2e^{|h||p|}$.
At derivative order zero use
$|e^u-1-u|\leq|u|^2e^{|u|}$; at order one use
$|e^u-1|\leq|u|e^{|u|}$; at higher orders differentiate the exponential,
which contributes at least two factors of $h$. These inequalities follow
directly from its power series. Choose $|h|<c/2$. Leibniz's rule and the
exponential decay of $L_z$ show that the remainder is $O(|h|^2)$ in
each Schwartz seminorm. The derivative map is continuous from the
finite-dimensional parameter space into $\mathscr S$, since multiplication
by each frequency coordinate is continuous. Composition with the continuous
linear functional $w$ proves complex differentiability of $H_n$;
in finite complex dimension this gives holomorphy.

For fixed $y$, put $g(p)=a_y(p)=L_{iy}(p)$. Modulation becomes translation
under Fourier transformation:
\[
 \widehat{L_{x+iy}}(u)=\widehat g(u-x).
\]
For $s_{k,j}(h)=\sup_u |u|^k\|D^j h(u)\|$, the triangle inequality
gives the explicit translation bound
\[
 s_{k,j}(h(\,\cdot-x))
 \leq2^k\big(s_{k,j}(h)+|x|^k s_{0,j}(h)\big).
\]
A continuous linear functional on Schwartz space is bounded by finitely
many such seminorms. Fourier inversion and the displayed bound therefore
bound $|w(L_{x+iy})|$ by a polynomial in $|x|$. Translation is continuous
in Schwartz space: its derivatives converge uniformly on compact sets,
and the same bound with one extra decay weight controls the tails.
Thus the integrand in \eqref{eq:tube-integral} is continuous and absolutely
integrable for every Schwartz $f$.

First let $f$ be compactly supported. Fourier inversion, pointwise in $p$,
gives
\[
 a_y(p)\check f(p)
  =(2\pi)^{-D}\int f(x)L_{x+iy}(p)\,dx.
\]
This integral can be passed through $w$ in the Schwartz topology.
For an explicit justification take Riemann sums in a box containing the
support of $f$. The integrand is continuous into Schwartz space and has
compact parameter support. For each seminorm the error is bounded by
the box volume times its uniform modulus of continuity in that seminorm.
The sums therefore converge in every seminorm to the displayed function;
continuity of $w$ gives \eqref{eq:tube-integral}. For general $f$, use
$f_R=\kappa(x/R)f$ with $0\leq\kappa\leq1$ and $\kappa=1$ near zero.
Then $f_R\to f$ in Schwartz space, while polynomial growth gives an
integrable majorant for $|H_n(v;x+iy)f_R(x)|$ independent of $R$.
Continuity on the left and dominated convergence on the right finish
the proof.
\end{proof}

\begin{proposition}[Identification with the ordered-time function]
\label{prop:tube-identification}
For strictly increasing scalar heights $y_1<\cdots<y_n$,
\[
 H_n(v;X+i d(y))=F_n(t+iy,x;v),
 \qquad X=((t_1,x_1),\ldots,(t_n,x_n)).
\]
For any $Y\in C_n$, the integrals of $H_n(v;X+iY')f(X)$ converge to
$W_Y(f)$ as $Y'\to Y$ through $C_n^\circ$.
\end{proposition}
\begin{proof}
At fixed strict scalar heights, $\psi_y$ is a tempered distribution by
the uniform real-position bound on $F_n$. Let $h$ be a compactly supported
smooth frequency test. Since $\chi=1$ near $K_n$ and $w$ is supported
there, the exponential testing identity gives
\[
 w(a_{d(y)}h)=w(e^{-d(y)\cdot p}h)
            =\psi_y(\widehat h).
\]
Both sides are continuous in $h$ on Schwartz space. Compact-test density
therefore extends this equality to every $h\in\mathscr S$.
Apply inverse Fourier transformation: $W_{d(y)}=\psi_y$.
Equation~\eqref{eq:tube-integral} now identifies the integrals of the
two proposed continuous functions against every Schwartz test. They agree
pointwise: a nonzero difference would have, after a constant phase change,
strictly positive real part on a ball and would be detected by a
nonnegative smooth bump there. The final boundary statement follows from
\eqref{eq:tube-integral} and continuity of $W_{Y'}(f)$ on $C_n$.
\end{proof}

\begin{proposition}[Lorentz covariance in the tube]\label{prop:tube-covariance}
For every proper orthochronous Lorentz matrix $\Lambda$,
\[
 H_n(v;z)=H_n(\rho(\Lambda)v;\Lambda z)
          \qquad(z\in\mathcal T_n).
\]
The action is simultaneous on all slots.
\end{proposition}
\begin{proof}
The matrix $\Lambda$ preserves strict forward gaps, so $\Lambda z$ is
in the tube. Its dual $A=\Lambda^{-T}$ preserves $K_n$, because it is
again proper orthochronous and preserves total momentum and suffix cones.
The Fourier form of Theorem~\ref{thm:lorentz-covariance}, obtained by
the change of variables of determinant one, is
\[
 \widehat W_n(\rho(\Lambda)v)(h\circ A^{-1})
       =\widehat W_n(v)(h).
\]
The two Schwartz functions $L_z$ and $L_{\Lambda z}\circ A$ agree
near $K_n$. Indeed their exponential factors agree because
$(Ap)\cdot(\Lambda z)=p\cdot z$. At each $p\in K_n$, both cutoffs
are one on the neighborhood
$B(p,1)\cap A^{-1}B(Ap,1)$. Their difference therefore vanishes near
the support of $w$. Apply the Fourier covariance identity to
$h=L_{\Lambda z}\circ A$, discard this vanishing difference, and multiply
by $(2\pi)^{-D}$. This is exactly the claimed equality.
\end{proof}
